% M/M/1 queue: the geometric stationary distribution of the number in
% system at two utilisations, and the mean number in system blowing up
% as utilisation approaches one.
\begin{figure}[htbp]
\centering
\begin{tikzpicture}[font=\small]

% ---- Left panel: geometric stationary law at two utilisations ----
\begin{axis}[
  name=geom,
  width=7.0cm, height=5.6cm,
  ybar, bar width=5pt,
  xlabel={number in system $k$},
  ylabel={$\Pr(N=k)=(1-\rho)\rho^{k}$},
  xmin=-0.6, xmax=12.6, ymin=0, ymax=0.62,
  xtick={0,2,4,6,8,10,12},
  ytick={0,0.2,0.4,0.6},
  tick label style={font=\scriptsize},
  label style={font=\scriptsize},
  title style={font=\small\bfseries},
  title={Stationary law},
  legend style={font=\scriptsize, at={(0.98,0.98)}, anchor=north east,
    draw=enggray!50},
]
  % rho = 0.5 : (1-0.5)0.5^k
  \addplot[fill=engblue!55, draw=engblue] coordinates {
    (0,0.5) (1,0.25) (2,0.125) (3,0.0625) (4,0.03125)
    (5,0.0156) (6,0.0078) (7,0.0039)
  };
  \addlegendentry{$\rho=0.5$}
  % rho = 0.8 : (0.2)0.8^k
  \addplot[fill=engred!45, draw=engred] coordinates {
    (0,0.2) (1,0.16) (2,0.128) (3,0.1024) (4,0.0819)
    (5,0.0655) (6,0.0524) (7,0.0419) (8,0.0336) (9,0.0268)
    (10,0.0215) (11,0.0172) (12,0.0137)
  };
  \addlegendentry{$\rho=0.8$}
\end{axis}

% ---- Right panel: mean number in system vs utilisation ----
\begin{axis}[
  at={($(geom.east)+(1.5cm,0)$)}, anchor=west,
  width=7.0cm, height=5.6cm,
  xlabel={utilisation $\rho=\lambda/\mu$},
  ylabel={mean in system $\mathrm{E}[N]=\dfrac{\rho}{1-\rho}$},
  xmin=0, xmax=1.0, ymin=0, ymax=20,
  xtick={0,0.25,0.5,0.75,1.0},
  ytick={0,5,10,15,20},
  grid=both, grid style={enggray!25},
  tick label style={font=\scriptsize},
  label style={font=\scriptsize},
  title style={font=\small\bfseries},
  title={Congestion near saturation},
]
  % E[N] = rho/(1-rho); clamp domain below 1 to avoid divergence
  \addplot[thick, engblue, domain=0:0.952, samples=120] {x/(1-x)};
  % markers at rho = 0.8 and 0.9
  \addplot[only marks, mark=*, engred, mark size=2pt]
    coordinates {(0.8,4.0) (0.9,9.0)};
  \draw[dashed, engred] (axis cs:0.8,0) -- (axis cs:0.8,4.0);
  \draw[dashed, engred] (axis cs:0.9,0) -- (axis cs:0.9,9.0);
  \node[anchor=south east, font=\scriptsize, engred] at (axis cs:0.80,4.4)
    {$\rho=0.8,\ \bar N=4$};
  \node[anchor=east, font=\scriptsize, engred] at (axis cs:0.88,11.0)
    {$\rho=0.9,\ \bar N=9$};
  % asymptote
  \draw[dotted, enggray] (axis cs:1.0,0) -- (axis cs:1.0,20);
\end{axis}

\end{tikzpicture}
\caption[The M/M/1 queue: geometric law and congestion]{The M/M/1 queue,
a single server with Poisson arrivals at rate $\lambda$ and exponential
service at rate $\mu$, has utilisation $\rho=\lambda/\mu$ and a geometric
stationary distribution $\Pr(N=k)=(1-\rho)\rho^{k}$ for the number in
system. Left: the law at $\rho=0.5$ decays fast, while at $\rho=0.8$ it
has a long tail, so long queues are common at high load. Right: the mean
number in system $\mathrm{E}[N]=\rho/(1-\rho)$ climbs slowly at low load and then
diverges as $\rho\to1$. Raising utilisation from $0.8$ to $0.9$, a
$12.5\%$ increase in load, more than doubles the mean occupancy
$\mathrm{E}[N]$ from $4$ to $9$, the congestion non-linearity that makes a
queue run near capacity fragile. A heavy-tailed service time makes matters worse still through the
Pollaczek-Khinchine correction. Source: computed by the authors from the
M/M/1 stationary law, \cite[ch.~6]{text:wasserman2004}.}
\label{fig:vol02:ch13:mm1-queue}
\end{figure}
